% Conversão de chapter-06-text-h-0109; origem RELATORIO_ATUAL.odt.
\hypertarget{__RefHeading___Toc4835_3118176633}{}
\refstepcounter{subsection}
\subsection*{6.2.2 Banco de dados utilizado}
\label{ages-iv:text-h-0109}
\addcontentsline{toc}{subsection}{6.2.2 Banco de dados utilizado}

% SOURCE text-p-1339
O sistema utiliza o banco de dados relacional PostgreSQL\space\cite{ages-iv-postgresql}, hospedado em ambiente gerenciado por meio da plataforma Supabase\space\cite{ages-iv-supabase}. A modelagem foi estruturada para contemplar as entidades centrais do projeto, como usuários, publicações (Recordas), relações de seguidores, preferências musicais e registros de moderação, garantindo integridade referencial, consistência transacional e segurança no controle de acesso aos dados. Além do armazenamento operacional do feed e das interações sociais, a estrutura foi projetada para viabilizar consultas analíticas de histórico e agregações de métricas musicais dos usuários, bem como rotinas de moderação e auditoria do painel administrativo.\par

% SOURCE text-a-0148 (âncora ODF vazia; destino preservado sem texto visível)
\href{https://tools.ages.pucrs.br/2026-2/2jk-4jk/recorda/wiki/-/wikis/Banco\%20de\%20Dados}{}

% FIGURE draw-image-0025 / CAPTION text-p-1344; imagem e legenda preservam o destino ODF.
\memorialfigurelinked{imagens/original/asset-0019.png}{\textbf{Figura 25 –} Diagrama Banco de Dados - AGES IV}{ages-iv:draw-image-0025}{Fonte: pendente de identificação}{0.92}{https://tools.ages.pucrs.br/2026-2/2jk-4jk/recorda/wiki/-/wikis/Banco\%20de\%20Dados}

% SOURCE text-p-1343 / text-a-0149
Mais detalhes podem ser encontrados aqui: \href{https://tools.ages.pucrs.br/2026-2/2jk-4jk/recorda/wiki/-/wikis/Banco\%20de\%20Dados}{Banco de Dados}\par
